\section{Discussion}
\label{sec:discussion}

\subsection{Diagnostic Utility Is a Property of the Decision Setting}
\label{sec:mismatch}

The same diagnostic decisions lead to different comparisons with always-graph when the candidate models change. The GCN and GAT restrictions make this dependence concrete: a rule that loses against the full portfolio can have lower mean regret against either restricted portfolio. Prior work characterizes graph properties and model performance~\citep{luan2023whendo,platonov2023characterizing,zheng2024trihom}; our evaluation connects such statistics to the cost of selecting a specified model portfolio. This makes the reference policy and candidate set part of the diagnostic's evaluation target.

Available improvement helps explain the contrast. Always-graph has only 0.256 percentage points of regret in the full portfolio, leaving little absolute improvement for any selector using these outcomes. The GCN and GAT portfolios leave more headroom. Absolute regret and improvement over a reference remain separate questions: GPR-GNN gives lower absolute Combined regret than GCN or GAT but favors always-graph in the paired comparison. Each restricted portfolio also defines its own oracle. The observed reversals therefore concern these model-selection settings, with favorable averages concentrated in Cornell and Wisconsin.

\subsection{Fallback and Training Determine the Measured Cost}

An abstaining diagnostic becomes an operational policy only after specifying what happens on abstention. Here MLP fallback turns Combined's full-set action into a single homophily threshold. The MLP validation threshold changes coverage, while fallback determines the selected predictions on uncovered units. The loss decomposition shows why covered-set regret alone gives an incomplete assessment: most of the full-set loss occurs after abstention. Reporting coverage, covered-set loss, and the full-set contribution of fallback makes this distinction explicit.

Training choices also change the comparison being predicted. The raw-feature experiment recovers substantial MLP accuracy on Roman-empire and reduces both tested graph--MLP gaps. The validation-only affine-input diagnostics provide further evidence that a shared preprocessing transform does not ensure an adequately trained feature-only baseline. These observations support inspecting baseline optimization before interpreting a graph-selection loss as a property of graph structure. The degree-preserving intervention separately demonstrates useful edge organization for GCN on the three citation networks; its graphs and architectures do not identify the mechanism of the largest diagnostic losses.

\subsection{Scope of the Evidence}

The benchmark covers 11 public datasets, one stratified split design, ten seeds, and a four-trial grid per architecture. Its simple rules test transparent decision policies; published diagnostics with richer features or learned selectors are a different comparison. Thresholds were fixed after legacy aggregate inspection and before the primary records. Subsequent portfolio restrictions, calibration, and influence checks are post-hoc. In particular, all three favorable subset averages disappear when Cornell and Wisconsin are jointly excluded, so transport to independent datasets remains unresolved.

The full graph action receives 24 trials against four for MLP. Single-architecture restrictions match trial counts while changing candidate availability, and neither analysis equalizes time or FLOPs or evaluates a comparably expanded MLP search. The two-dataset preprocessing experiment does not establish full-benchmark robustness to alternative inputs. Chameleon and Squirrel use original versions with known duplicate-node concerns; their exclusion sensitivity describes the remaining datasets, not retraining on filtered graphs.

Dataset-level inference has a concrete resolution limit. Fixed policy agreement on four datasets restricts the key comparison to at most seven nonzero differences, making rejection unattainable in the chosen Holm family. The descriptive loss comparison supplies the operational conclusion; non-rejection supplies no evidence of equivalence. More seeds on the same datasets would not remove those structural zeros. A new evaluation would need to plan dataset composition and inferential resolution together.

Finally, reconstruction of decisions from retained outcomes and repeatability of model training are separate properties. The bounded LINKX audit records default-CUDA variability and identical histories within its specified deterministic repeat groups. It characterizes the tested configurations rather than certifying the original benchmark's training. The released snapshot supports record verification; the private predecessor history and large checkpoint tensors remain in the author archive.

The next discriminating comparisons concern richer published diagnostics and stronger feature-only training under a specified budget. Their value would be to test whether the observed decision costs persist when the selector or its baseline changes. The current evidence already supports reporting the candidate portfolio, baseline training, fallback, and available headroom whenever diagnostic utility is evaluated.
